\documentclass[12pt]{article}
\usepackage{amsmath,amssymb}
\usepackage{booktabs}
\usepackage{geometry}
\usepackage{graphicx}
\usepackage{float}

\geometry{margin=1in}

\title{\textbf{Romberg Integration and Comparison with Simpson's Rule}}
\author{
Course: CSE261 -- Numerical Methods\\
Group Name: \textbf{Perseus}
}
\date{}

\begin{document}
\maketitle

%------------------------------------------------
\section{Introduction}

Let the function be defined as
\[
f(x) = \sin x
\]

The objective is to numerically evaluate the definite integral
\[
I = \int_a^b f(x)\,dx
\]
where
\[
a = 0 \quad \text{and} \quad b = 1
\]

The exact value of the integral is
\[
I = \int_0^1 \sin x\,dx = 1 - \cos(1) \approx 0.4596976941
\]

In this report, the integral is first approximated using the Composite Trapezoidal Rule.
Romberg Integration is then applied to accelerate the convergence of the Trapezoidal Rule
using Richardson extrapolation. Finally, the numerical results obtained using Romberg
Integration are compared with Simpson’s 1/3 Rule in terms of accuracy and convergence behavior.

%------------------------------------------------
\section{Composite Trapezoidal Rule}

The Composite Trapezoidal Rule approximates a definite integral by dividing the interval
$[a,b]$ into $n$ equal subintervals.

Let,
\[
a = \text{lower limit of integration}
\]
\[
b = \text{upper limit of integration}
\]
\[
n = \text{number of subintervals}
\]
\[
h = \frac{b-a}{n} \quad \text{(step size)}
\]

The nodal points are defined as
\[
x_i = a + ih, \quad i = 0,1,2,\ldots,n
\]
and
\[
y_i = f(x_i)
\]

The Composite Trapezoidal Rule is given by
\[
\int_a^b f(x)\,dx \approx T_n
\]
where
\[
T_n = \frac{h}{2}
\left[
(y_0 + y_n) + 2(y_1 + y_2 + \cdots + y_{n-1})
\right]
\]

\subsection*{Application}

For
\[
f(x) = \sin x, \quad [a,b] = [0,1]
\]
we take
\[
n = 4 \Rightarrow h = \frac{1-0}{4} = 0.25
\]

\begin{table}[H]
\centering
\begin{tabular}{c c}
\toprule
$x$ & $f(x)=\sin x$ \\
\midrule
0.00 & 0.0000000000 \\
0.25 & 0.2474039593 \\
0.50 & 0.4794255386 \\
0.75 & 0.6816387600 \\
1.00 & 0.8414709848 \\
\bottomrule
\end{tabular}
\caption{Tabulated values of $f(x)$}
\end{table}

\[
T_4 = \frac{0.25}{2}
\left[(0 + 0.8414709848) + 2(0.2474039593 + 0.4794255386 + 0.68163876)\right]
\]

\[
T_4 = 0.4573009376
\]

\[
\text{Absolute Error} = |0.4596976941 - 0.4573009376| = 0.0023967571
\]

%------------------------------------------------
\section{Romberg Integration}

Romberg Integration improves the Composite Trapezoidal Rule by applying Richardson
extrapolation to eliminate lower-order error terms.

Let
\[
R(k,0) = \text{Trapezoidal approximation using } 2^k \text{ subintervals}
\]

The recursive Romberg formula is
\[
R(k,j) = R(k,j-1) +
\frac{R(k,j-1) - R(k-1,j-1)}{4^j - 1}
\]

\subsection*{Romberg Table}

\begin{table}[H]
\centering
\begin{tabular}{c c c c}
\toprule
$k \backslash j$ & 0 & 1 & 2 \\
\midrule
0 & 0.4500805155 & & \\
1 & 0.4573009376 & 0.4597077450 & \\
2 & 0.4593451120 & 0.4596974570 & 0.4596976941 \\
\bottomrule
\end{tabular}
\caption{Romberg Integration Table}
\end{table}

The most accurate approximation is obtained from the last diagonal element:
\[
R(2,2) = 0.4596976941
\]

\[
\text{Absolute Error} \approx 0
\]

%------------------------------------------------
\section{Simpson's 1/3 Rule}

Simpson’s 1/3 Rule approximates the integral by fitting quadratic polynomials
over pairs of subintervals.

Let
\[
n = \text{even number of subintervals}
\]
\[
h = \frac{b-a}{n}
\]

The Simpson’s 1/3 Rule formula is
\[
\int_a^b f(x)\,dx \approx
\frac{h}{3}
\left[
(y_0 + y_n) + 4(y_1 + y_3 + \cdots + y_{n-1})
+ 2(y_2 + y_4 + \cdots + y_{n-2})
\right]
\]

For $n=4$ and $h=0.25$,
\[
S = \frac{0.25}{3}
\left[(0 + 0.8414709848) + 4(0.2474039593 + 0.68163876)
+ 2(0.4794255386)\right]
\]

\[
S = 0.4597077448
\]

\[
\text{Absolute Error} = 0.0000100507
\]

%------------------------------------------------
\section{Error and Convergence Analysis}

The Composite Trapezoidal Rule has an error order of $O(h^2)$, whereas Simpson’s Rule
has an error order of $O(h^4)$. Romberg Integration further accelerates convergence
by systematically eliminating lower-order error terms through Richardson extrapolation.

%------------------------------------------------
\section{Error Decay Comparison}

To illustrate convergence acceleration, the absolute error is plotted against the
refinement level for Romberg Integration and Simpson’s 1/3 Rule.

\begin{figure}[H]
\centering
\includegraphics[width=0.75\textwidth]{error_decay.png}
\caption{Error decay comparison between Romberg Integration and Simpson's Rule}
\end{figure}

The figure clearly shows that the error in Romberg Integration decreases much faster
than that of Simpson’s Rule, confirming convergence acceleration.

%------------------------------------------------
\section{Comparison of Results}

\begin{table}[H]
\centering
\begin{tabular}{c c c}
\toprule
Method & Approximate Value & Absolute Error \\
\midrule
Composite Trapezoidal Rule & 0.4573009376 & 0.0023967571 \\
Simpson’s 1/3 Rule & 0.4597077448 & 0.0000100507 \\
Romberg Integration & 0.4596976941 & $\approx 0$ \\
\bottomrule
\end{tabular}
\caption{Comparison of numerical integration methods}
\end{table}

%------------------------------------------------
\section{Conclusion}

The numerical experiments and error decay plots confirm that Romberg Integration
provides superior accuracy and faster convergence compared to both the Composite
Trapezoidal Rule and Simpson’s 1/3 Rule. Richardson extrapolation effectively
accelerates convergence, making Romberg Integration a powerful numerical
integration technique.

\end{document}
